\section{Method}
\label{sec:method}

SRP separates the sparse basis learned over readout rows from the scalar score
explained at a hidden state. We define targets first because they fix the
meaning of support and opposition.

\subsection{Scalar Readout Targets}
\label{sec:problem-setup}

The readout SAE in \Cref{sec:method-srp} can reconstruct the LM head and produce
target-independent readout-feature profiles. A local account adds one choice:
the scalar quantity to explain. We call any linear scalar choice a readout
target. Throughout, ``readout'' means the final unembedding or LM head
\(\WU\in\R^{|\mathcal{V}|\times d}\). For hidden state \(h\in\R^d\), row \(t\)
is \(w_t^\top\), so the full logit vector is \(\WU h\).

Formally, let \(\alpha\in\R^{|\mathcal{V}|}\) be a vector of vocabulary-row
coefficients and define
\[
    q_\alpha = \sum_{t\in\mathcal{V}} \alpha_t w_t,
    \qquad
    s_\alpha(h)=h^\top q_\alpha .
\]
The familiar logit-lens coordinate is the token-logit special case
\[
    \ell_A(h)=h^\top w_A,
\]
where \(\alpha_A=1\) and all other coefficients are zero. SRP can account for
token logits directly, but contrasts often reduce sensitivity to shared
output-embedding and tokenizer row structure
\citep{press2017outputembedding,inan2017tying,yang2018softmaxbottleneck,bostrom2020bpe}.

Target choice fixes the question: support and opposition are defined relative to
a specified scalar readout target. In the source-verification example,
\(w_{\texttt{verify}}\) asks what raises the \texttt{verify} logit;
\(w_{\texttt{verify}}-\bar{w}_{\mathcal{V}}\) asks what raises it above an
average vocabulary row; \(w_{\texttt{verify}}-w_{\texttt{assume}}\) asks why it
beats a specific alternative; and a local margin compares it with several nearby
competitors. These linear targets answer distinct readout questions.

The remaining target families are linear contrasts. A pairwise logit difference,
\[
    m_{A,B}(h)=h^\top(w_A-w_B),
\]
asks why the readout supports token \(A\) over token \(B\); shared offsets cancel
and the sign encodes the favored token. For answer, action, or abstention
behaviors spread across several tokens, we use group contrasts. For any nonempty
token set \(\mathcal{A}\subseteq\mathcal{V}\), write
\[
    \bar{w}_{\mathcal{A}}
    =
    \frac{1}{|\mathcal{A}|}\sum_{a\in\mathcal{A}}w_a .
\]
Then the group contrast is
\[
    q_{\mathcal{A},\mathcal{B}}
    =
    \bar{w}_{\mathcal{A}}-\bar{w}_{\mathcal{B}} .
\]
This keeps the target linear while reducing dependence on one tokenizer row. The
vocabulary-mean row
\(\bar{w}_{\mathcal{V}}=|\mathcal{V}|^{-1}\sum_{t\in\mathcal{V}}w_t\) gives a
reference contrast. We also use margins to nearby competitors,
\(w_u-\sum_{r\in\mathcal{R}}\pi_r w_r\), with fixed competitor set
\(\mathcal{R}\) and weights \(\pi\). The main target families are token logits,
reference contrasts, logit differences, group contrasts, and top-competitor
margins.

The scalar \(s_\alpha(h)\) is the unit of a local score account. Fixing it before
combining the target with the SRP basis keeps token preferences, reference
contrasts, and decision margins distinct.
